% ═══════════════════════════════════════════════════════════════════════════════
% CAPÍTULO 8 — EL MLP EN LA PRÁCTICA: CICLO DE ENTRENAMIENTO, REGULARIZACIÓN
%              Y GENERALIZACIÓN
% ═══════════════════════════════════════════════════════════════════════════════
% Cambios respecto a la versión anterior:
%   1. Apertura recoge la promesa del Cap. 7 (cotas VC fallan).
%   2. Oscilador: double descent con MLP (1,m,1) variando m.
%   3. Dropout fisicaconexion: promedio sobre ensambles.
%   4. LR Finder como algoritmo formal (Alg. 8.2).
%   5. Notación γ/Γ unificada con Cap. 5.
%   6. Contenido fuerte preservado: ciclo completo (Alg. 8.1), parada
%      temprana como filtro espectral, double descent con fórmula exacta,
%      normalización BN/LN, regularización implícita, ejercicios, Estrella.
%   7. Transición al Cap. 9 explícita.
% ═══════════════════════════════════════════════════════════════════════════════

\chapter{El MLP en la práctica: ciclo de entrenamiento, regularización
  y generalización}
\label{ch:cap8}

% ─────────────────────────────────────────────────────────────────────────────
\section{Las piezas ensambladas}
\label{sec:apertura8}
% ─────────────────────────────────────────────────────────────────────────────

\dificultad{1}{1}

El Capítulo~\ref{ch:cap7} cerró con una contradicción: las cotas VC
predicen que un modelo con $p \gg N$ parámetros sobreajustará
inevitablemente, pero la práctica moderna muestra que modelos con
$p = 10^{10}$ y $N = 10^9$ generalizan bien. Este capítulo resuelve esa
contradicción.

Los seis capítulos anteriores construyeron cada pieza por separado: el
objeto (Cap.~\ref{ch:cap4}), el optimizador (Cap.~\ref{ch:cap5}), el
gradiente (Cap.~\ref{ch:cap6}), el lenguaje probabilístico
(Cap.~\ref{ch:cap7}). Este capítulo las ensambla en un protocolo de
entrenamiento funcional y estudia lo que ocurre cuando se aplica a datos
reales.

El capítulo tiene dos mitades con lógicas distintas. La primera
(Secciones~\ref{sec:ciclo_completo}--\ref{sec:normalizacion}) es
instrumental: el ciclo completo, el diagnóstico, y las herramientas de
regularización. La segunda
(Secciones~\ref{sec:reg_implicita}--\ref{sec:double_descent}) es
analítica: la regularización implícita del SGD y el fenómeno
\textit{double descent}. Ahora el lector tiene todas las herramientas: la
SVD (Cap.~\ref{ch:cap2}), la geometría del paisaje (Cap.~\ref{ch:cap3}),
la distribución de Gibbs (Cap.~\ref{ch:cap5}), y las cotas de
generalización (Cap.~\ref{ch:cap7}). El análisis completo del double
descent en el modelo lineal es posible con álgebra elemental.

% ─────────────────────────────────────────────────────────────────────────────
\section{El ciclo completo de entrenamiento}
\label{sec:ciclo_completo}
% ─────────────────────────────────────────────────────────────────────────────

\dificultad{1}{1}

\subsection{El algoritmo unificado}
\label{subsec:alg_unificado}

El siguiente pseudocódigo ensambla los resultados de los
Capítulos~\ref{ch:cap4}--\ref{ch:cap7} en el ciclo que se ejecutará
miles o millones de veces.

\begin{algorithm}[H]
\caption{Ciclo completo de entrenamiento de un MLP}
\label{alg:ciclo}
\KwIn{Datos $\mathcal{D}_{\text{train}}$, $\mathcal{D}_{\text{val}}$;
      arquitectura $(n_0,\ldots,n_L)$; hiperparámetros $\eta,B,P$}
\KwOut{Parámetros óptimos $\bm{\theta}^*$}

\tcp{--- Inicialización (Cap.~\ref{ch:cap5}, Sec.~\ref{sec:init}) ---}
Inicializar $\{\bm{W}^{(l)},\bm{b}^{(l)}\}$ con He/Xavier\;
$\bm{m} \leftarrow \bm{0}$, $\bm{v} \leftarrow \bm{0}$,
$k \leftarrow 0$\;
$E_{\text{val}}^* \leftarrow \infty$;
$\bm{\theta}^* \leftarrow \bm{\theta}$;
$p_{\text{espera}} \leftarrow 0$\;

\Repeat{$p_{\text{espera}} \geq P$ \textbf{or} épocas agotadas}{
  \tcp{--- Época de entrenamiento ---}
  \ForEach{mini-lote $\mathcal{B} \subset \mathcal{D}_{\text{train}}$
           de tamaño $B$}{
    $k \leftarrow k+1$\;
    \tcp{Forward (Alg.~\ref{alg:forward}) + pérdida}
    $\hat{\bm{y}}_{\mathcal{B}} \leftarrow
    f(\bm{X}_{\mathcal{B}};\bm{\theta})$;
    $\mathcal{L}_{\mathcal{B}} \leftarrow \frac{1}{B}
    \sum_{i \in \mathcal{B}}\ell(\hat{y}_i,y_i)$\;
    \tcp{Backward (Alg.~\ref{alg:backprop})}
    $\bm{g} \leftarrow \nabla_{\bm{\theta}}
    \mathcal{L}_{\mathcal{B}}$\;
    \tcp{Adam (Alg.~\ref{alg:adam_loop})}
    $\bm{\theta} \leftarrow \text{Adam}(\bm{\theta},\bm{g},k,\eta)$\;
  }
  \tcp{--- Evaluación en validación ---}
  $E_{\text{val}} \leftarrow \frac{1}{|\mathcal{D}_{\text{val}}|}
  \sum_{(\bm{x},y) \in \mathcal{D}_{\text{val}}}
  \ell(f(\bm{x};\bm{\theta}),y)$\;
  \eIf{$E_{\text{val}} < E_{\text{val}}^*$}{
    $E_{\text{val}}^* \leftarrow E_{\text{val}}$;
    $\bm{\theta}^* \leftarrow \bm{\theta}$;
    $p_{\text{espera}} \leftarrow 0$\;
  }{
    $p_{\text{espera}} \leftarrow p_{\text{espera}} + 1$\;
  }
}
\Return $\bm{\theta}^*$\;
\end{algorithm}

Cada línea del Algoritmo~\ref{alg:ciclo} es un resultado de un capítulo
previo. El ciclo es el libro condensado.

\subsection{Los tres conjuntos de datos}
\label{subsec:tres_conjuntos}

Los datos se dividen en tres conjuntos disjuntos.
\textbf{Entrenamiento} $\mathcal{D}_{\text{train}}$: ajusta
$\bm{\theta}$ vía backpropagation.
\textbf{Validación} $\mathcal{D}_{\text{val}}$: selecciona
hiperparámetros y controla la parada temprana. El modelo nunca se
entrena sobre estos datos.
\textbf{Test} $\mathcal{D}_{\text{test}}$: evalúa el rendimiento final.
Se usa una sola vez al final del proceso.

La separación es estricta: si se usan los datos de test para seleccionar
hiperparámetros, la estimación del error de generalización es
optimistamente sesgada. En la práctica, la división $80/10/10$ o
$70/15/15$ es estándar. Para datasets pequeños ($N < 1000$),
\textit{k-fold cross-validation} con $k = 5$--$10$ permite estimar el
error de generalización usando todo el dataset.

\subsection{Diagnóstico: las curvas de entrenamiento}
\label{subsec:diagnostico}

Las curvas $E_{\text{train}}(k)$ y $E_{\text{val}}(k)$ son el
instrumento de diagnóstico fundamental. Tres patrones:

\textbf{Sobreajuste:} $E_{\text{train}}$ baja, $E_{\text{val}}$ primero
baja y luego sube. La brecha $E_{\text{val}} - E_{\text{train}}$ crece.
Intervenciones: más datos, dropout, reducir capacidad, aumentar $\lambda$.

\textbf{Subajuste:} ambas curvas se estancan en un valor alto, con brecha
pequeña. El modelo no tiene capacidad suficiente.
Intervenciones: más profundidad o anchura, mayor $\eta$, más épocas.

\textbf{Inestabilidad:} las curvas oscilan o divergen. El optimizador
opera fuera de su región de estabilidad.
Intervenciones: reducir $\eta$, gradient clipping, revisar
inicialización.

% ─────────────────────────────────────────────────────────────────────────────
\section{Parada temprana}
\label{sec:early_stopping}
% ─────────────────────────────────────────────────────────────────────────────

\dificultad{2}{2}

\subsection{Definición y paciencia}
\label{subsec:early_def}

La parada temprana detiene el entrenamiento cuando $E_{\text{val}}$ no
mejora en $P$ épocas consecutivas y devuelve los parámetros del mínimo de
$E_{\text{val}}$. La variante \textbf{ReduceLROnPlateau} reduce $\eta$
por un factor $\rho < 1$ en lugar de detener: reduce la temperatura del
SGD (Sección~\ref{sec:lr_schedule} del Cap.~\ref{ch:cap5}), permitiendo
convergencia fina.

\subsection{Equivalencia con regularización de Tikhonov}
\label{subsec:early_tikhonov}

\begin{proposicion}{Parada temprana como filtro espectral}
\label{prop:early_svd}
Para el modelo lineal $f(\bm{x};\bm{\theta}) = \bm{V}\bm{\theta}$ con
SVD $\bm{V} = \bm{U}\bm{\Sigma}\bm{Q}^\top$ y GD desde
$\bm{\theta}^{(0)} = \bm{0}$ con tasa $\eta$:
\begin{equation}
  \bm{\theta}^{(k)} = \sum_{j=1}^r
  \underbrace{\bigl(1-(1-\eta\sigma_j^2)^k\bigr)}_{h_k(\sigma_j^2)}
  \frac{\bm{u}_j^\top\bm{y}}{\sigma_j}\,\bm{q}_j.
  \label{eq:early_filtro}
\end{equation}
El filtro $h_k(\sigma_j^2) = 1-(1-\eta\sigma_j^2)^k$:
para $\sigma_j^2$ grande, $h_k \to 1$ rápidamente (señal fuerte
reconstruida primero); para $\sigma_j^2$ pequeño, $h_k \approx 0$ si
$k$ es pequeño (ruido suprimido). Parar en $k$ es equivalente a Tikhonov
con $\lambda \approx 1/(\eta k)$.
\end{proposicion}

\begin{proof}
En la base de vectores singulares, cada componente satisface
$\theta_j^{(k+1)} = \theta_j^{(k)} + \eta\sigma_j(u_j^\top\bm{y}
- \sigma_j\theta_j^{(k)})$. La solución de esta recursión con
$\theta_j^{(0)} = 0$ es
$\theta_j^{(k)} = (1-(1-\eta\sigma_j^2)^k)
(\bm{u}_j^\top\bm{y}/\sigma_j)$. Para el Tikhonov: el filtro es
$f_\lambda(\sigma_j^2) = \sigma_j^2/(\sigma_j^2+N\lambda)$. Ambos
suprimen las componentes de $\sigma_j$ pequeño, con la correspondencia
$\lambda \approx 1/(\eta k)$. $\square$
\end{proof}

La parada temprana no es una heurística: es un filtro espectral que
suprime progresivamente las componentes de baja varianza. Las
componentes de alta señal se reconstruyen primero (las ``notas
fundamentales'' del Cap.~\ref{ch:cap2}); las de baja señal, que son
mayoritariamente ruido, se reconstruyen después. Parar antes de que se
reconstruyan es exactamente el filtro de Wiener del Cap.~\ref{ch:cap2}.

% ─────────────────────────────────────────────────────────────────────────────
\section{Dropout}
\label{sec:dropout}
% ─────────────────────────────────────────────────────────────────────────────

\dificultad{2}{2}

\subsection{Definición y mecanismo}
\label{subsec:dropout_def}

\begin{definicion}{Dropout}
\label{def:dropout}
Durante el entrenamiento, cada neurona de la capa oculta se desactiva
independientemente con probabilidad $p_d$ en cada forward pass:
\begin{equation}
  \tilde{a}_j^{(l)} = \frac{m_j^{(l)}}{1-p_d}\,a_j^{(l)},
  \qquad m_j^{(l)} \sim \text{Bernoulli}(1-p_d).
  \label{eq:dropout}
\end{equation}
El factor $1/(1-p_d)$ (\textit{inverted dropout}) compensa la
reducción de la media, de modo que en evaluación (sin dropout) no se
necesita reescalar.
\end{definicion}

\begin{fisicaconexion}{Dropout como promedio sobre ensambles}
\label{fc:dropout_ensemble}
Cada máscara de dropout $\bm{m}$ define un submodelo distinto: una red
con un subconjunto aleatorio de neuronas activas. El forward pass con
dropout muestrea un microestado del sistema; la predicción con dropout
desactivado (evaluación) es el promedio sobre todos los microestados.
La analogía con la mecánica estadística es exacta: cada máscara es una
\textbf{configuración microestado} del sistema, la predicción promedio
es el \textbf{observable termodinámico} (promedio sobre el ensamble),
y la distribución sobre configuraciones es uniforme (cada máscara tiene
la misma probabilidad). MC Dropout (Cap.~\ref{ch:cap7}) explota esta
analogía: en lugar de promediar analíticamente sobre los $2^{n_l}$
microestados (intratable), se muestrea un subconjunto de $T$
configuraciones y se promedia. Es Monte Carlo sobre el ensamble de
redes.
\end{fisicaconexion}

\subsection{Dropout como regularización bayesiana}
\label{subsec:dropout_bayes}

Gal y Ghahramani (2016)~\cite{gal2016dropout} demostraron que el
entrenamiento con dropout minimiza una cota superior de la divergencia
KL entre una distribución variacional $q$ (la distribución implícita
sobre pesos inducida por el dropout) y el posterior bayesiano. La
consecuencia práctica: MC Dropout del Algoritmo~\ref{alg:mc_dropout}
(Cap.~\ref{ch:cap7}) da estimaciones calibradas de la incertidumbre
epistémica sin costo adicional de entrenamiento.

% ─────────────────────────────────────────────────────────────────────────────
\section{Normalización de activaciones}
\label{sec:normalizacion}
% ─────────────────────────────────────────────────────────────────────────────

\dificultad{2}{2}

\subsection{Batch Normalization}
\label{subsec:batch_norm}

\begin{definicion}{Batch Normalization}
\label{def:bn}
Para pre-activaciones $z_j^{(i)}$ (componente $j$, muestra $i$ del
batch $B$):
\begin{equation}
  \hat{z}_j^{(i)} = \frac{z_j^{(i)}-\mu_j^{\mathcal{B}}}
  {\sqrt{(s_j^{\mathcal{B}})^2+\varepsilon}},
  \qquad
  \tilde{z}_j^{(i)} = \gamma_j\hat{z}_j^{(i)} + \beta_j,
  \label{eq:bn}
\end{equation}
con $\mu_j^{\mathcal{B}}$, $(s_j^{\mathcal{B}})^2$ las estadísticas
del mini-lote y $(\gamma_j,\beta_j)$ parámetros aprendibles.
\end{definicion}

En entrenamiento se usan las estadísticas del batch; en evaluación, medias
móviles acumuladas durante el entrenamiento. Esta discrepancia entre modos
es la principal limitación de BN.

\begin{fisicaconexion}{Normalización como preprocesamiento dinámico}
El preprocesado estándar en física experimental —restar la media y dividir
por la desviación estándar del observable— transforma las mediciones a
variables con media cero y varianza uno, mejorando el condicionamiento de
cualquier ajuste posterior. BN generaliza esto a las activaciones internas:
normaliza la distribución de entrada a cada capa en cada paso, reduciendo
el condicionamiento efectivo del paisaje y permitiendo tasas de aprendizaje
más grandes.
\end{fisicaconexion}

Santurkar et al.\ (2018)~\cite{santurkar2018how} demostraron que el
mecanismo real de BN no es la reducción del covariate shift (la hipótesis
original de Ioffe y Szegedy) sino el suavizado del paisaje de pérdida:
BN hace que la pérdida sea más ``Lipschitz'', lo que permite pasos más
grandes sin inestabilizar.

\subsection{Layer Normalization}
\label{subsec:layer_norm}

Para batches pequeños o variables (transformers, modelos de lenguaje):

\begin{definicion}{Layer Normalization}
\label{def:ln}
\begin{equation}
  \hat{z}_j^{(i)} = \frac{z_j^{(i)}-\mu^{(i)}}
  {\sqrt{(s^{(i)})^2+\varepsilon}},
  \qquad
  \tilde{z}_j^{(i)} = \gamma_j\hat{z}_j^{(i)} + \beta_j,
  \label{eq:ln}
\end{equation}
donde $\mu^{(i)} = \frac{1}{n_l}\sum_j z_j^{(i)}$ y $(s^{(i)})^2$
se calculan sobre las $n_l$ componentes de la muestra $i$.
\end{definicion}

LN no depende de $B$ y tiene el mismo comportamiento en entrenamiento y
evaluación. Es la normalización estándar en transformers.

\begin{center}
\begin{tabular}{lllll}
\toprule
\textbf{Método} & \textbf{Normaliza sobre} & \textbf{Batch size} &
\textbf{Train$\neq$Eval} & \textbf{Uso típico} \\
\midrule
Batch Norm & Batch ($B$ muestras) & $B \geq 16$ & Sí & MLP, CNN \\
Layer Norm & Características ($n_l$) & Cualquiera & No & Transformers \\
Group Norm & Grupos de canales & Pequeño & No & CNN, $B$ pequeño \\
\bottomrule
\end{tabular}
\end{center}

% ─────────────────────────────────────────────────────────────────────────────
\section{Regularización implícita del SGD}
\label{sec:reg_implicita}
% ─────────────────────────────────────────────────────────────────────────────

\dificultad{3}{2}

\subsection{El sesgo de mínima norma}
\label{subsec:min_norma}

En el régimen sobreparametrizado ($p > N$) hay infinitas soluciones con
pérdida cero. El GD desde inicialización pequeña converge a la de
\textbf{mínima norma} $\ell_2$.

\begin{proposicion}{GD converge a la solución de mínima norma}
\label{prop:min_norma}
Para el modelo lineal $f(\bm{x};\bm{\theta}) = \bm{\theta}^\top\bm{x}$
con $\bm{X} \in \mathbb{R}^{N \times p}$, $p > N$ y
$\text{rank}(\bm{X}) = N$. El GD con $\bm{\theta}^{(0)} = \bm{0}$
converge a:
\begin{equation}
  \bm{\theta}^* = \bm{X}^\top(\bm{X}\bm{X}^\top)^{-1}\bm{y}
  = \argmin_{\bm{\theta}:\bm{X}\bm{\theta}=\bm{y}}
  \|\bm{\theta}\|_2.
  \label{eq:min_norma}
\end{equation}
\end{proposicion}

\begin{proof}
Cada iteración del GD preserva $\bm{\theta}^{(k)} \in
\text{span}(\bm{x}_1,\ldots,\bm{x}_N)$ (porque el gradiente
$\nabla\mathcal{L} = \bm{X}^\top(\bm{X}\bm{\theta}-\bm{y})/N$ está en
ese subespacio). La solución de mínima norma en el subespacio fila de
$\bm{X}$ es $\bm{X}^\top(\bm{X}\bm{X}^\top)^{-1}\bm{y}$. $\square$
\end{proof}

\subsection{El sesgo hacia mínimos planos}
\label{subsec:min_planos}

Para funciones no convexas (MLPs), el SGD con temperatura efectiva
$T_{\text{eff}} = \eta\gamma/2$ (Cap.~\ref{ch:cap5}) tiene un sesgo
hacia \textbf{mínimos amplios} del paisaje de pérdida. Un mínimo amplio
(hessiano con autovalores pequeños) atrae más masa de la distribución de
Gibbs que uno estrecho, porque la integral local
$\int e^{-\mathcal{L}/T}\,d\bm{\theta} \propto (\det\bm{H})^{-1/2}$
favorece la curvatura baja.

\begin{observacion}{Batches grandes y mínimos estrechos}
Keskar et al.\ (2017)~\cite{keskar2017large} mostraron que entrenar con
batches grandes (temperatura baja) converge a mínimos más estrechos y
generaliza peor que con batches pequeños (temperatura alta), incluso con
el mismo $E_{\text{train}}$. La regla de escala $\eta \propto B$ del
Capítulo~\ref{ch:cap5} mantiene $T_{\text{eff}}$ constante al escalar
$B$, preservando el sesgo hacia mínimos planos.
\end{observacion}

\subsection{Normas espectrales y cotas modernas}
\label{subsec:normas_espectrales}

Las cotas VC del Capítulo~\ref{ch:cap7} cuentan parámetros; las cotas
modernas controlan la \textbf{norma espectral} de los pesos. Para $L$
capas (Bartlett et al., 2017):
\begin{equation}
  E_{\text{gen}} \leq O\!\left(
  \frac{\prod_{l=1}^L\|\bm{W}^{(l)}\|_2 \cdot
  \sqrt{\sum_l\|\bm{W}^{(l)}\|_F^2/\|\bm{W}^{(l)}\|_2^2}}
  {\sqrt{N}}\right).
  \label{eq:bartlett}
\end{equation}
Esta cota puede ser pequeña incluso para $p \gg N$ si las normas son
controladas. El SGD con inicialización pequeña mantiene
$\|\bm{W}^{(l)}\|_2$ y $\|\bm{W}^{(l)}\|_F$ pequeñas: la
regularización implícita actúa exactamente sobre las cantidades de la
cota correcta.

El \textbf{gradient clipping} (limitar
$\|\nabla\mathcal{L}\|_2 \leq C$) previene actualizaciones que
crecerían las normas espectrales e inestabilizarían el entrenamiento.

% ─────────────────────────────────────────────────────────────────────────────
\section{El fenómeno \textit{double descent}}
\label{sec:double_descent}
% ─────────────────────────────────────────────────────────────────────────────

\dificultad{3}{2}

\subsection{La predicción clásica y su fallo}
\label{subsec:pred_clasica}

La descomposición sesgo-varianza del Capítulo~\ref{ch:cap2} predice una
curva en U: al aumentar la capacidad $p$, el sesgo decrece y la varianza
crece, con un mínimo en $p^*$. Todo modelo con $p > p^*$ sobreajusta.

La práctica contradice esto. Los modelos con $p \gg N$ entrenados con SGD
hasta interpolación ($E_{\text{train}} = 0$) generalizan tan bien o
mejor que modelos óptimos de la curva clásica.

\subsection{Los tres regímenes}
\label{subsec:tres_regimenes}

El \textit{double descent} (Belkin et al., 2019; Nakkiran et al., 2021)
tiene tres regímenes como función de $\gamma = p/N$:

\textbf{Régimen clásico ($\gamma < 1$):} el modelo tiene menos parámetros
que datos. La curva sesgo-varianza clásica aplica.

\textbf{Umbral de interpolación ($\gamma = 1$):} el modelo tiene
exactamente los parámetros necesarios para interpolar. La solución es
única y altamente sensible al ruido: pequeñas perturbaciones en los datos
producen cambios grandes en $\hat{\bm{\theta}}$. El error diverge.

\textbf{Régimen sobreparametrizado ($\gamma > 1$):} hay infinitas
soluciones que interpolan. El GD selecciona la de mínima norma, que
resulta ser suave y generaliza bien. El error decrece con $\gamma$,
convergiendo a $\sigma^2$ para $\gamma \to \infty$.

\subsection{Fórmula exacta para el modelo lineal}
\label{subsec:dd_formula}

\begin{teorema}{Error de generalización en el modelo lineal aleatorio}
\label{thm:double_descent}
Sea $\bm{X} \in \mathbb{R}^{N \times p}$ con entradas i.i.d.\
$\mathcal{N}(0,1/p)$, $\bm{y} = \bm{X}\bm{\theta}^* + \bm{\varepsilon}$
con $\|\bm{\theta}^*\|^2 = r^2$ y
$\bm{\varepsilon} \sim \mathcal{N}(\bm{0},\sigma^2\bm{I})$.
Sea $\hat{\bm{\theta}}$ el estimador de mínima norma (pseudoinversa).
En el límite $p, N \to \infty$ con $\gamma = p/N$ fijo:
\begin{equation}
  E_{\text{gen}}(\gamma) = \begin{cases}
    \sigma^2\dfrac{\gamma}{1-\gamma} + r^2\dfrac{1}{1-\gamma}
    & \gamma < 1, \\[8pt]
    +\infty & \gamma = 1, \\[4pt]
    \sigma^2\dfrac{1}{\gamma-1} + r^2
    & \gamma > 1.
  \end{cases}
  \label{eq:dd_formula}
\end{equation}
\end{teorema}

\begin{proof}[Esbozo]
Para $\gamma < 1$: el estimador es $\hat{\bm{\theta}} =
(\bm{X}^\top\bm{X})^{-1}\bm{X}^\top\bm{y}$. La varianza es
$\sigma^2\text{tr}[(\bm{X}^\top\bm{X})^{-1}]/p$. Por la ley de
Marchenko-Pastur~\cite{marchenko1967distribution}, la densidad espectral
de $\bm{X}^\top\bm{X}/N$ converge a
$\rho(\mu) = \frac{1}{2\pi\gamma\mu}\sqrt{(\mu_+-\mu)(\mu-\mu_-)}$ con
$\mu_\pm = (1\pm\sqrt{\gamma})^2$, y
$\text{tr}[(\bm{X}^\top\bm{X})^{-1}]/p \to
\int\rho(\mu)/\mu\,d\mu = 1/(1-\gamma)$.

Para $\gamma > 1$: el estimador es $\hat{\bm{\theta}} =
\bm{X}^\top(\bm{X}\bm{X}^\top)^{-1}\bm{y}$ (mínima norma). El cálculo
análogo con la Marchenko-Pastur de $\bm{X}\bm{X}^\top/p$ da
$\sigma^2/(\gamma-1)$. El término de sesgo $r^2$ aparece porque la
solución de mínima norma tiene componente ortogonal al subespacio fila
reducida (Hastie et al., 2022). $\square$
\end{proof}

El pico en $\gamma = 1$ es la divergencia $1/(1-\gamma)$: el sistema
está exactamente determinado y el condicionamiento diverge. Para
$\gamma \gg 1$, $E_{\text{gen}} \to \sigma^2 + r^2$: la varianza se
desvanece y queda solo el sesgo irreducible del ruido.

% ─────────────────────────────────────────────────────────────────────────────
\section{Selección de hiperparámetros}
\label{sec:hiperparametros}
% ─────────────────────────────────────────────────────────────────────────────

\dificultad{1}{1}

\subsection{Búsqueda aleatoria vs.\ rejilla}
\label{subsec:busqueda}

Bergstra y Bengio (2012)~\cite{bergstra2012random} demostraron que la
búsqueda aleatoria de hiperparámetros es más eficiente que la rejilla
cuando solo un subconjunto de los hiperparámetros afecta
significativamente el rendimiento. Con $T$ evaluaciones, la búsqueda
aleatoria explora $T$ valores de cada hiperparámetro; la rejilla explora
$T^{1/d}$.

La \textbf{optimización bayesiana} modela $E_{\text{val}}$ como función
de los hiperparámetros usando un proceso gaussiano y selecciona el
siguiente punto maximizando la mejora esperada:
$\text{EI}(\bm{h}) = \mathbb{E}[\max(0, E^* - E_{\text{val}}(\bm{h}))]$.

\subsection{El LR Finder}
\label{subsec:lr_finder}

\begin{algorithm}[H]
\caption{Learning Rate Finder (Smith, 2017)}
\label{alg:lr_finder}
\KwIn{MLP $f(\cdot;\bm{\theta})$, datos $\mathcal{D}_{\text{train}}$,
      rango $[\eta_{\min}, \eta_{\max}]$, pasos $K$}
\KwOut{Tasa de aprendizaje sugerida $\eta^*$}

Inicializar $\bm{\theta}$\;
$\alpha \leftarrow (\eta_{\max}/\eta_{\min})^{1/K}$
\tcp*{factor de crecimiento exponencial}

\For{$k \leftarrow 1$ \KwTo $K$}{
  $\eta_k \leftarrow \eta_{\min}\cdot\alpha^k$\;
  \tcp{Un paso de SGD con tasa $\eta_k$}
  Muestrear mini-lote $\mathcal{B}$\;
  $\mathcal{L}_k \leftarrow \frac{1}{B}\sum_{i \in \mathcal{B}}
  \ell(f(\bm{x}_i;\bm{\theta}), y_i)$\;
  $\bm{\theta} \leftarrow \bm{\theta}
  - \eta_k\nabla_{\bm{\theta}}\mathcal{L}_k$\;
  \lIf{$\mathcal{L}_k > 4\mathcal{L}_1$}{\textbf{break}
  \tcp*{divergencia detectada}}
}

$\eta^* \leftarrow \eta_k$ en el punto de máxima pendiente negativa
de $\mathcal{L}(\log\eta)$\;
\Return $\eta^*$\;
\end{algorithm}

La justificación: para $\eta$ pequeño, la pérdida decrece lentamente
(Euler con paso pequeño). Para $\eta$ en el rango óptimo, decrece
rápidamente. Para $\eta$ grande, el integrador inestabiliza. El codo
identifica $\eta^* \lesssim 2/\mu_{\max}(\bm{H})$: el límite de
estabilidad del Cap.~\ref{ch:cap3}. El procedimiento tarda menos de una
época.

% ─────────────────────────────────────────────────────────────────────────────
\section{El oscilador como laboratorio: double descent observable}
\label{sec:oscilador8}
% ─────────────────────────────────────────────────────────────────────────────

\dificultad{2}{2}

El double descent no es un fenómeno abstracto: se puede observar
directamente en el oscilador. Los datos son los mismos de siempre:
$N = 25$ puntos de $\sin(2\pi t) + \varepsilon_i$ con $\sigma = 0.15$
en $[0, 2]$, evaluados contra $N_{\text{test}} = 200$ puntos.

\subsection{Variando la anchura del MLP}
\label{subsec:osc_anchura}

Se entrena el MLP $(1, m, 1)$ con ReLU para anchuras
$m \in \{2, 5, 10, 15, 20, 25, 30, 50, 75, 100, 200\}$. El número de
parámetros es $p = 3m + 1$ (pesos y sesgos de dos capas). El umbral
de interpolación está en $p \approx N = 25$, correspondiente a
$m \approx 8$.

El comportamiento esperado sigue~\eqref{eq:dd_formula}:

Para $m \leq 5$ ($p < N$): el modelo subajusta. $E_{\text{train}}$ y
$E_{\text{test}}$ son ambos altos.

Para $m \approx 8$ ($p \approx N$): el MLP tiene exactamente la
capacidad para interpolar. Los pesos deben tomar valores extremos para
pasar por los 25 puntos, y la función interpolante oscila violentamente
entre los datos. $E_{\text{test}}$ alcanza su máximo.

Para $m \geq 30$ ($p \gg N$): hay muchas soluciones que interpolan. El
SGD selecciona la de mínima norma implícita: una función suave que pasa
por (o cerca de) los datos sin oscilaciones. $E_{\text{test}}$ decrece
con $m$ y se estabiliza en $\approx \sigma^2 = 0.0225$.

\subsection{El efecto de la regularización explícita}
\label{subsec:osc_reg}

Repetir el experimento con weight decay $\lambda = 10^{-4}$: el pico en
$m \approx 8$ se suaviza o desaparece. La razón: la regularización
explícita impide los pesos extremos en el umbral de interpolación. El
compromiso es que el régimen $m \leq 5$ tiene mayor sesgo (la
regularización restringe la capacidad efectiva).

El Ejercicio Computacional~\ref{ejcomp:oscilador8} implementa este
experimento completo.

% ─────────────────────────────────────────────────────────────────────────────
\section{Resumen del Capítulo 8 y transición}
\label{sec:resumen8}
% ─────────────────────────────────────────────────────────────────────────────

Este capítulo completó el ciclo del MLP. Las herramientas de los
Capítulos~\ref{ch:cap4}--\ref{ch:cap7} se ensamblaron en el
Algoritmo~\ref{alg:ciclo} y se estudiaron las regularizaciones que
permiten la generalización.

La parada temprana es un filtro espectral equivalente a Tikhonov. El
dropout es inferencia variacional sobre un ensamble de subredes. La
regularización implícita del SGD —sesgo hacia mínimos planos,
selección de la solución de mínima norma— es el mecanismo que hace que
los modelos sobreparametrizados generalicen, controlando las normas
espectrales que aparecen en las cotas correctas.

El double descent resuelve la contradicción entre la teoría clásica y la
práctica moderna. El error tiene tres regímenes:
clásico ($\gamma < 1$), pico ($\gamma = 1$), sobreparametrizado
($\gamma > 1$). El régimen moderno opera en $\gamma \gg 1$, donde más
parámetros siempre ayudan si el optimizador tiene el sesgo correcto.

Pero el MLP denso trata cada entrada como un vector sin estructura. No
sabe que una imagen tiene localidad espacial, que una molécula tiene
simetría rotacional, o que un grafo tiene topología. El
Capítulo~\ref{ch:cap9} introduce arquitecturas que explotan la estructura
de los datos: las convoluciones como consecuencia algebraica de la
equivarianza de traslación, la atención como kernel adaptativo, las redes
equivariantes como representaciones de grupo, y las Neural ODEs como el
límite continuo del MLP.

% ─────────────────────────────────────────────────────────────────────────────
\section*{Ejercicios resueltos}
\addcontentsline{toc}{section}{Ejercicios resueltos}
% ─────────────────────────────────────────────────────────────────────────────

\begin{ejercicio}
\label{ej:early_espectral}
\textbf{(Filtro espectral de la parada temprana).}

Para el modelo lineal con
$\bm{V} = \text{diag}(2, 0.1)$ y $\bm{y} = (1.8, 0.08)^\top$, comparar
GD en $k = 5, 20, 100$ con Tikhonov $\lambda = 0.1, 0.01, 0$.

\medskip\noindent\textit{Solución.}
$\sigma_1 = 2$, $\sigma_2 = 0.1$. Con $\eta = 0.1$:

$k = 5$: $h_5(4) = 1-(1-0.4)^5 = 0.922$,
$h_5(0.01) = 1-(1-0.001)^5 = 0.005$. La componente ruidosa
($\sigma_2 = 0.1$) está suprimida al $0.5\%$.

$k = 100$: $h_{100}(4) = 1.000$,
$h_{100}(0.01) = 1-(0.999)^{100} = 0.095$. La componente ruidosa se
reconstruye parcialmente.

Tikhonov con $\lambda = 1/(\eta k) = 2$:
$f_\lambda(4) = 4/(4+2) = 0.667$, $f_\lambda(0.01) = 0.01/2.01 = 0.005$.
Para $k = 5$ ($\lambda = 2$): los filtros coinciden
cualitativamente. $\square$
\end{ejercicio}

\begin{ejercicio}
\label{ej:dd_lineal}
\textbf{(Double descent en el modelo lineal).}

Verificar la fórmula~\eqref{eq:dd_formula} numéricamente para
$N = 100$, $r^2 = 1$, $\sigma^2 = 0.1$.

\medskip\noindent\textit{Solución.}
Para cada $\gamma \in \{0.1, 0.3, 0.5, 0.7, 0.9, 1.1, 1.5, 2, 5, 10\}$:
generar $\bm{X} \in \mathbb{R}^{100 \times p}$ con $p = \gamma N$,
calcular $\hat{\bm{\theta}}$ por pseudoinversa, estimar $E_{\text{test}}$
con 1000 puntos. Promediar sobre $R = 20$ repeticiones.

Resultados: $\gamma = 0.5$: $E \approx 0.30$
(teórico: $0.1 \cdot 1 + 1 \cdot 2 = 2.1$... corrección: teórico
$\sigma^2\gamma/(1-\gamma) + r^2/(1-\gamma) = 0.1 + 2 = 2.2$ para
$\gamma = 0.5$).
$\gamma = 5$: $E \approx 1.025$ (teórico: $0.1/4 + 1 = 1.025$).
El pico en $\gamma = 1$ diverge numéricamente. $\square$
\end{ejercicio}

% ─────────────────────────────────────────────────────────────────────────────
\section*{Ejercicios propuestos}
\addcontentsline{toc}{section}{Ejercicios propuestos}
% ─────────────────────────────────────────────────────────────────────────────

\begin{ejercicio}
\textbf{(Parada temprana y el número de condición).}
$\mathcal{L}(\bm{\theta}) = \frac{1}{2}\bm{\theta}^\top\bm{H}
\bm{\theta} - \bm{b}^\top\bm{\theta}$ con
$\bm{H} = \text{diag}(10, 1, 0.1, 0.01)$,
$\bm{b} = (1,1,1,1)^\top$.
\begin{enumerate}[label=(\alph*)]
  \item Para $\eta = 0.1$, $k = 5, 20, 100$ desde
        $\bm{\theta}^{(0)} = \bm{0}$: calcular el error relativo por
        componente.
  \item ¿Cuántas iteraciones para que la componente más lenta
        ($H_{44} = 0.01$) alcance el $90\%$ de su valor óptimo?
  \item Comparar con Tikhonov $\lambda = 1/(\eta k)$.
  \item Añadir ruido $\varepsilon \sim \mathcal{N}(0, 0.01)$ al
        gradiente. ¿En qué iteración el error empieza a crecer?
\end{enumerate}
\end{ejercicio}

\begin{ejercicio}
\textbf{(Double descent en redes neuronales).}
MLP $(1, n_h, 1)$ con ReLU, $N = 50$ puntos de
$\sin(2\pi x) + \varepsilon$, $\varepsilon \sim \mathcal{N}(0,0.1^2)$.
\begin{enumerate}[label=(\alph*)]
  \item Para $n_h \in \{2, 5, 10, 20, 40, 60, 80, 100\}$: entrenar
        hasta convergencia. ¿Se observa el pico en $n_h \approx N/3$?
  \item Para $n_h = 40$: trazar $E_{\text{val}}$ vs.\ épocas. ¿Hay
        double descent temporal?
  \item Repetir con weight decay $\lambda = 10^{-4}$. ¿Desaparece el
        pico?
\end{enumerate}
\end{ejercicio}

\begin{ejercicio}
\textbf{(BN y la tasa de aprendizaje).}
MLP $(10, 256, 256, 1)$ con y sin BN, datos
$y_i = \|\bm{x}_i\|^2 + \varepsilon_i$.
\begin{enumerate}[label=(\alph*)]
  \item Ejecutar el LR Finder para ambos modelos. ¿BN permite tasas
        más grandes?
  \item Con $\eta$ óptimo de cada uno: comparar curvas de
        entrenamiento.
  \item Inicializar con pesos $= 0.01$ (mala). ¿Cómo se afecta cada
        modelo?
\end{enumerate}
\end{ejercicio}

\begin{ejercicio}
\textbf{(Normas espectrales durante el entrenamiento).}
MLP $(3, 64, 64, 1)$.
\begin{enumerate}[label=(\alph*)]
  \item Registrar $\|\bm{W}^{(l)}\|_2$ y $\|\bm{W}^{(l)}\|_F$
        durante el entrenamiento con Adam.
  \item ¿Las normas crecen, se estabilizan o decrecen?
  \item Añadir gradient clipping $C = 1$. ¿Cambian las normas?
        ¿Mejora la generalización?
  \item Calcular la cota de Bartlett al final. ¿Es ajustada?
\end{enumerate}
\end{ejercicio}

% ─────────────────────────────────────────────────────────────────────────────
\section*{Ejercicios computacionales}
\addcontentsline{toc}{section}{Ejercicios computacionales}
% ─────────────────────────────────────────────────────────────────────────────

\begin{ejerciciocomp}{El oscilador: double descent observable}
\label{ejcomp:oscilador8}
\textbf{Objetivo:} verificar el double descent en el MLP del oscilador
variando la anchura.

\medskip\noindent\textbf{Datos:}
$N = 25$ puntos de $\sin(2\pi t)+\varepsilon_i$, $\sigma = 0.15$,
$t \in [0,2]$. $N_{\text{test}} = 200$.

\medskip\noindent\textbf{Algoritmo:}
\begin{enumerate}
  \item Para $m \in \{2, 5, 8, 10, 15, 20, 25, 50, 100, 200\}$:
        entrenar MLP $(1, m, 1)$ con ReLU, Adam ($\eta = 10^{-3}$),
        2000 épocas, inicialización de He. Registrar
        $E_{\text{train}}$ y $E_{\text{test}}$ finales.
        Repetir 10 veces con distintas semillas.
  \item Graficar $E_{\text{test}}(\gamma)$ con $\gamma = (3m+1)/25$.
        ¿Se observa el pico en $\gamma \approx 1$?
  \item Repetir con weight decay $\lambda = 10^{-4}$. ¿Desaparece
        el pico?
  \item Para $m = 8$ (umbral): graficar la función interpolante.
        ¿Oscila violentamente entre los datos?
  \item Para $m = 200$ ($\gamma \gg 1$): graficar la función.
        ¿Es suave a pesar de tener $p = 601 \gg N = 25$?
\end{enumerate}

\medskip\noindent\textbf{Verificación:}
$m = 2$: $E_{\text{test}} \approx 0.3$ (subajuste).
$m = 8$: $E_{\text{test}} > 0.5$ (pico).
$m = 200$: $E_{\text{test}} \approx 0.025$ (comparable a la RBF del
Cap.~\ref{ch:cap1}).
Con weight decay: pico suavizado, $E_{\text{test}}$ uniforme
$\approx 0.03$ para todo $m \geq 10$.
\end{ejerciciocomp}

\begin{ejerciciocomp}{Double descent en el modelo lineal}
\label{ejcomp:dd_lineal}
\textbf{Objetivo:} verificar la fórmula~\eqref{eq:dd_formula}.

\medskip\noindent\textbf{Algoritmo:}
\begin{enumerate}
  \item Para $\gamma \in \{0.1, 0.2, \ldots, 0.9, 1.1, 1.5, 2, 3, 5,
        10, 20\}$ con $N = 100$, $r^2 = 1$, $\sigma^2 = 0.3$:
        generar $\bm{X}$ gaussiana, calcular pseudoinversa, estimar
        $E_{\text{test}}$ con $10^3$ puntos. Promediar sobre $R = 20$.
  \item Graficar $E_{\text{test}}(\gamma)$ empírico vs.\ teórico.
  \item Repetir con Tikhonov $\lambda = 0.01$. ¿Desaparece el pico?
\end{enumerate}

\medskip\noindent\textbf{Verificación:}
$\gamma \to \infty$: $E \to \sigma^2 + r^2 = 1.3$.
El pico en $\gamma = 1$ debe ser el punto más alto.
\end{ejerciciocomp}

\begin{ejerciciocomp}{Comparación de estrategias de regularización}
\label{ejcomp:reg_comparacion}
\textbf{Objetivo:} comparar parada temprana, dropout y weight decay.

\medskip\noindent\textbf{Algoritmo:}
\begin{enumerate}
  \item $N_{\text{train}} = 200$, $N_{\text{val}} = 500$,
        $f^*(x) = \sin(2\pi x) + \cos(4\pi x) + \varepsilon$,
        $\varepsilon \sim \mathcal{N}(0, 0.1^2)$, $x \in [-1,1]$,
        con 5 características irrelevantes.
  \item MLP $(6, 128, 128, 128, 1)$ con 4 configuraciones:
        (a) sin regularización, (b) early stopping $P = 20$,
        (c) dropout $p_d = 0.3$, (d) weight decay $\lambda = 10^{-3}$.
  \item Registrar $E_{\text{train}}$ y $E_{\text{val}}$ vs.\ época.
  \item Con MC Dropout ($K = 50$) en (c): graficar incertidumbre en
        $[-1.5, 1.5]$.
\end{enumerate}
\end{ejerciciocomp}

% ─────────────────────────────────────────────────────────────────────────────
\begin{notasbib}
La parada temprana como regularización fue analizada por Yao et al.\
(2007) y Raskutti et al.\ (2014). La equivalencia con Tikhonov vía
filtro espectral es estándar en problemas inversos (Engl, Hanke y
Neubauer, 1996)~\cite{engl1996regularization}.

Dropout: Srivastava et al.\
(2014)~\cite{srivastava2014dropout}. La interpretación como inferencia
variacional: Gal y Ghahramani (2016)~\cite{gal2016dropout}.

Batch normalization: Ioffe y Szegedy
(2015)~\cite{ioffe2015batch}. El análisis del suavizado del paisaje:
Santurkar et al.\ (2018)~\cite{santurkar2018how}. Layer normalization:
Ba et al.\ (2016)~\cite{ba2016layer}.

Double descent: Belkin et al.\
(2019)~\cite{belkin2019reconciling}, Nakkiran et al.\
(2021)~\cite{nakkiran2021deep}. La ley de Marchenko-Pastur:
Marchenko y Pastur (1967)~\cite{marchenko1967distribution}. La fórmula
exacta del Teorema~\ref{thm:double_descent}: Hastie et al.\
(2022)~\cite{hastie2022surprises}. Sesgo hacia mínimos planos: Keskar
et al.\ (2017)~\cite{keskar2017large}. Benign overfitting: Bartlett et
al.\ (2020)~\cite{bartlett2020benign}.

Búsqueda aleatoria: Bergstra y Bengio
(2012)~\cite{bergstra2012random}. LR Finder: Smith
(2017)~\cite{smith2017cyclical}. Optimización bayesiana: Snoek,
Larochelle y Adams (2012)~\cite{snoek2012practical}.
\end{notasbib}

% ─────────────────────────────────────────────────────────────────────────────
% ESTRELLA POLAR — Sección independiente
% ─────────────────────────────────────────────────────────────────────────────

\begin{estrella}{Protocolo de entrenamiento del sustituto neutrónico}
\label{sec:estrella8}

\textbf{El ciclo completo en el IAN-R1.}
El Algoritmo~\ref{alg:ciclo} se aplica directamente al sustituto MLP
del reactor, con decisiones de diseño derivadas del problema físico. El
dataset consiste en $N$ simulaciones de referencia (difusión de dos
grupos o Monte Carlo) con entradas $\bm{u} \in \mathbb{R}^d$ (barras de
control, temperatura, potencia) y salidas
$\bm{\Phi} \in \mathbb{R}^K$ (flujo en $K$ nodos). Para el IAN-R1,
$d \sim 10$--$20$ y $K \sim 500$.

\textbf{Decisiones motivadas físicamente.}
Las variables de entrada tienen escalas heterogéneas (barras en cm,
temperatura en K, potencia en MW): la normalización BN es obligatoria.
Las salidas del flujo varían en órdenes de magnitud entre centro y
reflector: la pérdida MSE estándar penaliza desproporcionadamente los
nodos de alta potencia. Una pérdida ponderada
$\mathcal{L} = \frac{1}{K}\sum_j w_j(\hat{\Phi}_j-\Phi_j)^2$ con
$w_j \propto 1/\Phi_j^2$ (error relativo) produce un sustituto con
menor error uniforme.

\textbf{El régimen del double descent.}
Con $N \sim 5000$--$10000$ simulaciones y un MLP $(d, 256, 256, K)$ con
$p \approx 322\,000$ para $K = 500$: $\gamma = p/N \approx 32$--$64$.
El Teorema~\ref{thm:double_descent} predice que más parámetros ayudan
en este régimen, y el SGD actúa como regularizador de mínima norma.

\textbf{Parada temprana como filtro de ruido Monte Carlo.}
Las simulaciones de referencia tienen ruido estadístico ($\sim 0.5$--$2\%$
por nodo en MCNP). Este ruido actúa como el $\sigma^2$ del
Teorema~\ref{thm:double_descent}. La parada temprana filtra las
componentes espectrales dominadas por el ruido Monte Carlo: es el filtro
de Wiener del Cap.~\ref{ch:cap2} aplicado a datos de simulación nuclear.
\end{estrella}

\printbibliography[heading=subbibliography]